\BourbakiUnnumberedChapter{参考文献}{参考文献}

\begin{enumerate}
\def\labelenumi{(\Roman{enumi})}
\tightlist
\item
  《欧几里得原本》，5卷，J. L. 海伯格编，莱比锡（托伊布纳），1883-88年。
\end{enumerate}

（I bis）T. L. 希思，《欧几里得几何原本》十三卷\ldots，3卷，剑桥，1908年。

\begin{enumerate}
\def\labelenumi{(\Roman{enumi})}
\setcounter{enumi}{1}
\tightlist
\item
  《丢番图全集》\ldots，2卷，P. 坦纳里编，莱比锡（托伊布纳），1893-95年。
\end{enumerate}

（II bis）T. L. 希思，《亚历山大的丢番图》，第2 \(^{nd}\) 版，剑桥，1910年。

\begin{enumerate}
\def\labelenumi{(\Roman{enumi})}
\setcounter{enumi}{2}
\item
  B. 达塔与A. N. 辛格，《印度数学史》，2卷，拉合尔（莫蒂拉尔·巴纳尔西·达斯出版社），1935-38年。
\item
  S. 斯蒂文，《数学著作》\ldots，A. 吉拉德编，莱顿（爱思唯尔出版社），1634 年，第 I 卷。
\item
  L. 欧拉：a)《无穷分析引论》（《全集》，(1) 系列，第IX卷，苏黎世-莱比锡-柏林（O. Füssli 与 B. G. Teubner 出版社），1945年，第384页）；b)《固体或刚体运动理论》（《全集》(2) 系列，第III卷，苏黎世-莱比锡-柏林（O. Füssli 与 B. G. Teubner 出版社），1948年，第200-201页）；c)《代数学完全入门》（《全集》(1) 系列，第I卷，莱比锡-柏林（Teubner 出版社），1911年，第422页）；d)《关于方程虚根的研究》（《全集》(1) 系列，第VI卷，莱比锡-柏林（Teubner 出版社），1921年，第78页）。
\item
  J.-L. 拉格朗日，《著作集》，巴黎（Gauthier-Villars 出版社），1867-1892年：a)《积分学中若干问题的解法》，第I卷，第520页；b)《关于交点运动长期方程的研究》，第VI卷，第655-666页；c)《算术研究》，第III卷，第695-795页；d)《关于方程虚根的形式》，第III卷，第479页；e)《不受任何加速力作用的任意物体旋转问题的新解法》，第III卷，第579-616页。
\item
  P. S. 拉普拉斯：a) 关于微分方程的特解及行星长期不等式的论文（《著作集》，第 VIII 卷，巴黎（戈蒂埃-维拉尔出版社），1891 年，第 325-366 页）；b) 关于行星和卫星长期不等式的论文（《著作集》，第 XI 卷，巴黎（戈蒂埃-维拉尔出版社），1895 年，第 49-92 页）。
\item
  C. F. 高斯，《著作集》，第 I 卷（哥廷根，1870 年），第 II 卷（同上，1876 年）及第 III 卷（同上，1876 年）。
\end{enumerate}

(VIII bis) 高斯关于整代数函数分解为一次或二次实因子的四个证明（奥斯特瓦尔德经典著作，第14号，莱比锡（托伊布纳），1904年）。

\begin{enumerate}
\def\labelenumi{(\Roman{enumi})}
\setcounter{enumi}{8}
\item
  N. H. 阿贝尔，《全集》，第 I 卷，西洛与李编，克里斯蒂安尼亚，1881 年。
\item
  A. L. 柯西：a)《无穷小计算在几何中的应用讲义》（《全集》第二辑，第五卷，巴黎（戈蒂埃-维拉尔出版社），1903年，第248页）；b)《关于用以确定行星长期不等式的方程》（《全集》第二辑，第九卷，巴黎（戈蒂埃-维拉尔出版社），1891年，第174页）。
\item
  P. G. 勒热纳-狄利克雷，《著作集》，第 I 卷，柏林（G. 赖默），1889 年，第 619-644 页。
\item
  E. 库默尔，《关于复数理论》，J. de Crelle，第 XLIII 卷（1847 年），第 319 页（《论文集》，第 I 卷，海德堡（Springer 出版社），1975 年，第 203 页）。
\item
  Ch. 埃尔米特，《著作集》，第 I 卷，巴黎（Gauthier-Villars 出版社），1905 年。
\item
  J. J. 西尔维斯特，《数学论文集》，第 I 卷，剑桥，1904 年：关于二阶直线与曲面接触的枚举，第 219 页（=《哲学杂志》，1851 年）。
\item
  W. R. 哈密顿，《四元数讲义》，都柏林，1853 年。
\item
  A. 凯莱，《数学论文集》，剑桥，1889-1898 年：关于矩阵理论的回忆录，第 II 卷，第 475-496 页（=《哲学汇刊》，1858 年）。
\item
  H. J. 史密斯，《数学论文集》，第 I 卷，牛津，1894 年；关于线性不定方程组与同余式组，第 367 页（《哲学汇刊》，1861 年）。
\item
  E. 谢林，《可复合算术形式的基本类》，哥廷根科学学会论文集，第 XIV 卷（1868-69 年），第 13 页。
\item
  K. 魏尔斯特拉斯，《数学著作》，第 II 卷，柏林（Mayer und Müller 出版社），1895 年：关于双线性与二次形式理论，第 19 页。
\item
  L. 克罗内克，《关于理想复数类数的若干性质的阐述》，柏林科学院月报（1870 年），第 881 页（=《著作集》，第 I 卷，莱比锡（Teubner 出版社），1895 年，第 273 页）。
\item
  C. 若尔当，《置换与代数方程论》。巴黎（Gauthier-Villars 出版社），1870 年，第 114-125 页。
\item
  G. 弗罗贝尼乌斯，《具有整数系数的线性形式理论》，《全集》，第 I 卷，海德堡（Springer 出版社），1968 年，第 482 页（= J. de Crelle，1879 年）。
\item
  G. 弗罗贝尼乌斯与 L. 施蒂克尔贝格，《关于可交换元素的群》，J. de Crelle，第 LXXXVI 卷（1879 年），第 217 页（= 弗罗贝尼乌斯，《全集》，第 I 卷，第 545 页）。
\item
  R. 戴德金，《数学全集》，第 II 卷，不伦瑞克（Vieweg 出版社），1932 年：关于数通过其最大公因子的分解，第 103 页。
\item
  \begin{enumerate}
  \def\labelenumii{\alph{enumii})}
  \tightlist
  \item
    E. 阿廷与 O. 施赖埃尔，《实域的代数构造》，汉堡大学数学讨论班论文集，第 V 卷（1927 年），第 83 页；b) E. 阿廷，《关于确定函数分解为平方和》（同上，第 100 页）；c) E. 阿廷与 O. 施赖埃尔，《实闭域的一种刻画》（同上，第 225 页）。
  \end{enumerate}
\end{enumerate}

